\section*{Acknowledgments}
The authors are thankful for discussions with J.~H.~Zhang and Y.~B.~Yang. This material was partially supported by the U.S. Department of Energy, Office of Science, Office of Nuclear Physics from DE-FG02-93ER-40762, DE-SC0011090 and DE-SC0012704, by the Laboratory Directed Research and Development (LDRD) funding of BNL under contract DE-EC0012704, by a grant from the National Science Foundation of China (No.~11405104), and within the framework of the TMD Topical Collaboration.
%
I.S. was also supported in part by the Simons Foundation through the Investigator grant 327942.
T.I. was also supported by JSPS KAKENHI Grant Numbers JP26400261, JP17H02906, and also MEXT as ``Priority Issue on Post-K computer'' (Elucidation of the Fundamental Laws and Evolution of the Universe) and JICFuS.


\appendix

\section{Fourier Transform}
\label{sec:ft}

To Fourier transform the bare pseudo-PDF in Eq.~(\ref{eq:1looppseudo}) into the bare quasi-PDF, we use the identity
\begin{align} \label{eq:ft}
&\int {d \zeta \over 2\pi}\, e^{i x \zeta} (z^2\mu^2)^{\epsilon} \Gamma(-\epsilon) 4^{-\epsilon}\\
= & \left({\mu^2\over p_z^2}\right)^\epsilon \int {d \zeta \over 2\pi}\, e^{i x \zeta} \int_0^\infty d\alpha\, \alpha^{-1-\epsilon}e^{-\alpha\zeta^2}4^{-\epsilon}\nn\\
=& \left({\mu^2\over p_z^2}\right)^\epsilon \int_0^\infty {d\alpha\over 2\sqrt{\pi}}\, \alpha^{-3/2-\epsilon}e^{-x^2/(4\alpha)}4^{-\epsilon}\nn\\
=& \left({\mu^2\over p_z^2}\right)^\epsilon {\Gamma(\epsilon+1/2)\over \sqrt{\pi}} {1\over |x|^{1+2\epsilon}}\,,\nn
\end{align}
which is true for $\epsilon<0$.

Now let us turn to a plus function $g(y)^{[0,1]}_{+(1)}$. The double Fourier transform
\begin{align}
& \int {d \zeta \over 2\pi}\, e^{i x \zeta} \int_0^1 dy\, e^{-iy\zeta} g(y)^{[0,1]}_{+(1)} (z^2\mu^2)^{\epsilon} \Gamma(-\epsilon) 4^{-\epsilon}\nn\\
= &(i\zeta)\int {d \zeta \over 2\pi}\, e^{i (x-1) \zeta} \int_0^1 dy \int_0^1 dt\ y\,g(1-y)\nn\\
&\times e^{ity\zeta} (z^2\mu^2)^{\epsilon} \Gamma(-\epsilon) 4^{-\epsilon} \nn\\
=& \left({\mu^2\over p_z^2}\right)^\epsilon {\Gamma(\epsilon+{1\over2})\over \sqrt{\pi}}{\partial\over \partial x} \int_0^1 dy \int_0^1 dt \, {y\,g(1-y) \over |x-1+ty|^{1+2\epsilon}}\,.
\end{align}
Since the integration over $y$ and $t$ leads to a piecewise function of $x$, the derivative will end up with a plus function as the discontinuity at $x=1$ gives $\delta(1-x)$.


\section{Quasi-PDF Calculation}
\label{sec:quasi}

Here we quote the pure dimensional regularization results obtained when we carry out the quasi-PDF calculation following Refs.~\cite{Xiong:2013bka,Stewart:2017tvs} by Fourier transforming from $z$ to $x p^z$ for the integrand to obtain $[\delta(p^z-xp^z)-\delta(k^z-xp^z)]$.
For the vertex and wavefunction renormalization graphs we obtain
\begin{widetext}

\begin{align}
	&\tilde{q}^{(1)}_\text{vertex}(x,p^z,\epsilon)
	+ \tilde{q}^{(1)}_\text{w.fn.}(x,p^z,\epsilon)\nn\\
	&={\alpha_sC_F\over 2\pi}\left({\mu^2\over p_z^2}\right)^\epsilon \left[ (1-\epsilon)\int_0^1 dy (1-y){1\over |x-y|^{1+2\epsilon}}{\Gamma(\epsilon+{1\over2})\over\sqrt{\pi}}- \delta(1-x){1\over2} \left({1\over\epsilon_{\mbox{\tiny UV}}} - {1\over \epsilon_{\mbox{\tiny IR}}}\right)\right]\,,
\end{align}
while for the sail diagram
\begin{align}
	&\tilde{q}^{(1)}_\text{sail}(x,p^z,\epsilon) \nn\\
	&= {\alpha_sC_F\over 2\pi}\left({4\pi\mu^2\over p_z^2}\right)^\epsilon\left[\int_0^1 dy {x+y\over 1-x}{1\over |x-y|^{1+2\epsilon}}{\Gamma(\epsilon+{1\over2})\over\sqrt{\pi}}-\delta(1-x)\int dx'\int_0^1 dy {x'+y\over 1-x'}{1\over |x'-y|^{1+2\epsilon}}{\Gamma(\epsilon+{1\over2})\over\sqrt{\pi}} \right]\nonumber\\
	&= {\alpha_sC_F\over 2\pi}\left({\mu^2\over p_z^2}\right)^\epsilon\left[\int_0^1 dy {x+y\over 1-x}{1\over |x-y|^{1+2\epsilon}}{\Gamma(\epsilon+{1\over2})\over\sqrt{\pi}}+\delta(1-x)\left({1\over\epsilon_{\mbox{\tiny UV}}} - {1\over \epsilon_{\mbox{\tiny IR}}}\right)\right.\nn\\
	&\qquad\left. -\delta(1-x)\int dx'\int_0^1 dy {1+y\over 1-x'}{1\over |x'-y|^{1+2\epsilon}}{\Gamma(\epsilon_{\mbox{\tiny IR}}+{1\over2})\over\sqrt{\pi}}\right]\,,
	\end{align}
and for the tadpole diagram
	\begin{align}
	&\tilde{q}^{(1)}_\text{tadpole}(x,p^z,\epsilon)\nn\\
	=& {\alpha_sC_F\over 2\pi}\left({4\pi\mu^2\over p_z^2}\right)^\epsilon \left[-{1\over 1-2\epsilon}{1\over |1-x|^{1+2\epsilon}}{\Gamma(\epsilon+{1\over2})\over\sqrt{\pi}} -\delta(1-x) {1\over 1-2\epsilon}{1\over |1-x|^{1+2\epsilon}}\int_{-\infty}^\infty dx' {1\over |1-x'|^{1+2\epsilon}}{\Gamma(\epsilon+{1\over2})\over\sqrt{\pi}}\right]\nonumber\\
	=& {\alpha_sC_F\over 2\pi}\left({4\pi\mu^2\over p_z^2}\right)^\epsilon \left[-{1\over 1-2\epsilon}{1\over |1-x|^{1+2\epsilon}}{\Gamma(\epsilon+{1\over2})\over\sqrt{\pi}} +\delta(1-x)\left({1\over\epsilon_{\mbox{\tiny UV}}} - {1\over \epsilon_{\mbox{\tiny IR}}}\right) \right]\,.
	\end{align}

\end{widetext}
After adding these expressions and carrying out the remaining integrations over $y$, we obtain the same result as in Eq.~(\ref{eq:1loopquasi}).


\section{$\epsilon$ Expansion and Plus Functions}
\label{sec:ep}

Since the support of the quasi-PDF ranges from $-\infty$ to $\infty$, its asymptotic behavior as $\sim1/|x|$ at $|x|\to\infty$ implies a UV divergence which can be regularized by dimensional regularization. Therefore, the $\epsilon$ expansion of the quasi-PDF should account for this feature.

In general, we need to expand
\beq
	{\theta(x)\over x^{1+\epsilon}} = {\theta(x)\theta(1-x)\over x^{1+\epsilon}} + {\theta(x-1)\over x^{1+\epsilon}}\,,
\eeq
and it is well known that
\begin{align} \label{eq:ident}
{\theta(x)\theta(1-x)\over x^{1+\epsilon}} = - {1\over\epsilon}\delta(x)+
\tilde{\mathcal{L}}_0(x) - \epsilon \tilde{\mathcal{L}}_1(x)+ O(\epsilon^2)\,,
\end{align}
where $\epsilon<0$.
Here we follow~\cite{Ligeti:2008ac} and the plus functions $\mathcal{L}_n(x)$ are defined as
\begin{align} \label{eq:plus1}
	\mathcal{L}_n(x)\equiv &\left[ {\theta(x)\ln^nx\over x}\right]_+
   \\
	= &\lim_{\beta\to0} \left[{\theta(x-\beta)\ln^n x\over x} + \delta(x-\beta) {\ln^{n+1}\beta \over n+1}\right]
 \,, \nn
\end{align}
and we let
\begin{align}
  \tilde{\mathcal{L}}_n(x) = \theta(1-x) \mathcal{L}_n(x) \,.
\end{align}
Note that  $\int_0^1 dx \, {\cal L}_n(x) =0$.

To define the expansion in the range $x\in [1,\infty]$ we simply map this interval to $[0,1]$ via $t=1/x$. Taking an arbitrary smooth test function $g(x)$ we have
\begin{align}
	&\int_1^\infty dx {1\over x^{1+\epsilon}} \: g(x)
    \\
	=& \int_0^1 dt {1\over t^{1-\epsilon}}  \:  g(1/t)
   \nn\\
	=& \int_0^1 dt \left[{1\over\epsilon}\delta(t)+ \mathcal{L}_0(t) + \epsilon \mathcal{L}_1(t)+ O(\epsilon^2)\right] \, g(1/t)
    \nn\\
	=& \int_1^\infty\!\! {dx\over x^2}\left[{1\over\epsilon}\delta^+\Bigl({1\over x}\Bigr)+ \mathcal{L}_0\Bigl({1\over x}\Bigr) + \epsilon \mathcal{L}_1\Bigl({1\over x}\Bigr)+ O(\epsilon^2)\right] g(x)
   \,.\nn
\end{align}
Here $\epsilon>0$ and the superscript $+$ on the $\delta^+$ function indicates that its argument should be positive. Therefore $\delta^+(1/x)$ has its support at $x=+\infty$, not $x=-\infty$.  Since $g$ is arbitrary we can identify
\begin{align}
  {\theta(x-1)\over x^{1+\epsilon}}
   &= {1\over\epsilon}{1\over x^2}\delta^+\Bigl({1\over x}\Bigr)
   + \frac{1}{x^2} \tilde{\mathcal{L}}_0\Bigl({1\over x}\Bigr)
   + \epsilon \frac{1}{x^2} \tilde{\mathcal{L}}_1\Bigl({1\over x}\Bigr)
  \nn\\
  &\quad + O(\epsilon^2)\,.
\end{align}
Combining the above results and denoting which $1/\epsilon$ poles are UV or IR divergences we have
\begin{align} \label{eq:expansion}
	{\theta(x)\over x^{1+\epsilon}}
 = & \left[-\frac{1}{\epsilon_{\mbox{\tiny IR}}}\delta(x) + \frac{1}{\epsilon_{\mbox{\tiny UV}}} {1\over  x^2}\delta^+\Bigl({1\over x}\Bigr)\right]
    \\
    & + \left({1\over x}\right)^{[0,1]}_{+(0)}+\left({1\over x}\right)^{[1,\infty]}_{+(\infty)}
  \nn\\
	& - \epsilon\left[\left({\ln x\over x}\right)^{[0,1]}_{+(0)}+\left({\ln x\over x}\right)^{[1,\infty]}_{+(\infty)}\right] + O(\epsilon^2)
   \,, \nn
\end{align}
where we have defined the distributions
\begin{align}
  \big(1/x\big)_{+(0)}^{[0,1]} &\equiv  \tilde{{\cal L}}_0(x)
  \,,  \quad
  \big(\ln x/x\big)_{+(0)}^{[0,1]} \equiv  \tilde{{\cal L}}_1(x)
  \,,
\end{align}
and
\begin{align}\label{eq:plusinfinity}
 \big(1/x\big)_{+(\infty)}^{[1,\infty]} &\equiv (1/x^2) \tilde{{\cal L}}_0(1/x)
   \, \\
\big(\ln x/x\big)_{+(\infty)}^{[1,\infty]} &\equiv -(1/x^2) \tilde{{\cal L}}_1(1/x)
   \,. \nn
\end{align}
Note that \eq{expansion} is consistent with the expected result that
\begin{align}
 \int_0^\infty \frac{dx}{x^{1+\epsilon}}
 = \frac{1}{\epsilon_{\mbox{\tiny UV}} }
  -\frac{1}{\epsilon_{\mbox{\tiny IR}} } \,.
\end{align}


\begingroup
\makeatletter
%merlin.mbs apsrev4-1.bst 2010-07-25 4.21a (PWD, AO, DPC) hacked
%Control: key (0)
%Control: author (8) initials jnrlst
%Control: editor formatted (1) identically to author
%Control: production of article title (-1) disabled
%Control: page (0) single
%Control: year (1) truncated
%Control: production of eprint (0) enabled
\begin{thebibliography}{69}%
\makeatletter
\providecommand \@ifxundefined [1]{%
 \@ifx{#1\undefined}
}%
\providecommand \@ifnum [1]{%
 \ifnum #1\expandafter \@firstoftwo
 \else \expandafter \@secondoftwo
 \fi
}%
\providecommand \@ifx [1]{%
 \ifx #1\expandafter \@firstoftwo
 \else \expandafter \@secondoftwo
 \fi
}%
\providecommand \natexlab [1]{#1}%
\providecommand \enquote  [1]{``#1''}%
\providecommand \bibnamefont  [1]{#1}%
\providecommand \bibfnamefont [1]{#1}%
\providecommand \citenamefont [1]{#1}%
\providecommand \href@noop [0]{\@secondoftwo}%
\providecommand \href [0]{\begingroup \@sanitize@url \@href}%
\providecommand \@href[1]{\@@startlink{#1}\@@href}%
\providecommand \@@href[1]{\endgroup#1\@@endlink}%
\providecommand \@sanitize@url [0]{\catcode `\\12\catcode `\$12\catcode
  `\&12\catcode `\#12\catcode `\^12\catcode `\_12\catcode `\%12\relax}%
\providecommand \@@startlink[1]{}%
\providecommand \@@endlink[0]{}%
\providecommand \url  [0]{\begingroup\@sanitize@url \@url }%
\providecommand \@url [1]{\endgroup\@href {#1}{\urlprefix }}%
\providecommand \urlprefix  [0]{URL }%
\providecommand \Eprint [0]{\href }%
\providecommand \doibase [0]{http://dx.doi.org/}%
\providecommand \selectlanguage [0]{\@gobble}%
\providecommand \bibinfo  [0]{\@secondoftwo}%
\providecommand \bibfield  [0]{\@secondoftwo}%
\providecommand \translation [1]{[#1]}%
\providecommand \BibitemOpen [0]{}%
\providecommand \bibitemStop [0]{}%
\providecommand \bibitemNoStop [0]{.\EOS\space}%
\providecommand \EOS [0]{\spacefactor3000\relax}%
\providecommand \BibitemShut  [1]{\csname bibitem#1\endcsname}%
\let\auto@bib@innerbib\@empty
%</preamble>
\bibitem [{\citenamefont {Collins}\ \emph {et~al.}(1989)\citenamefont
  {Collins}, \citenamefont {Soper},\ and\ \citenamefont
  {Sterman}}]{Collins:1989gx}%
  \BibitemOpen
  \bibfield  {author} {\bibinfo {author} {\bibfnamefont {J.~C.}\ \bibnamefont
  {Collins}}, \bibinfo {author} {\bibfnamefont {D.~E.}\ \bibnamefont {Soper}},
  \ and\ \bibinfo {author} {\bibfnamefont {G.~F.}\ \bibnamefont {Sterman}},\
  }\href@noop {} {\bibfield  {journal} {\bibinfo  {journal} {Adv. Ser. Direct.
  High Energy Phys.}\ }\textbf {\bibinfo {volume} {5}},\ \bibinfo {pages} {1}
  (\bibinfo {year} {1989})},\ \Eprint {http://arxiv.org/abs/hep-ph/0409313}
  {arXiv:hep-ph/0409313 [hep-ph]} \BibitemShut {NoStop}%
%%CITATION = HEP-PH/0409313;%%
\bibitem [{\citenamefont {Ball}\ \emph {et~al.}(2015)\citenamefont {Ball} \emph
  {et~al.}}]{Ball:2014uwa}%
  \BibitemOpen
  \bibfield  {author} {\bibinfo {author} {\bibfnamefont {R.~D.}\ \bibnamefont
  {Ball}} \emph {et~al.} (\bibinfo {collaboration} {NNPDF}),\ }\href@noop {}
  {\bibfield  {journal} {\bibinfo  {journal} {JHEP}\ }\textbf {\bibinfo
  {volume} {04}},\ \bibinfo {pages} {040} (\bibinfo {year} {2015})},\ \Eprint
  {http://arxiv.org/abs/1410.8849} {arXiv:1410.8849 [hep-ph]} \BibitemShut
  {NoStop}%
%%CITATION = ARXIV:1410.8849;%%
\bibitem [{\citenamefont {Dulat}\ \emph {et~al.}(2016)\citenamefont {Dulat},
  \citenamefont {Hou}, \citenamefont {Gao}, \citenamefont {Guzzi},
  \citenamefont {Huston}, \citenamefont {Nadolsky}, \citenamefont {Pumplin},
  \citenamefont {Schmidt}, \citenamefont {Stump},\ and\ \citenamefont
  {Yuan}}]{Dulat:2015mca}%
  \BibitemOpen
  \bibfield  {author} {\bibinfo {author} {\bibfnamefont {S.}~\bibnamefont
  {Dulat}}, \bibinfo {author} {\bibfnamefont {T.-J.}\ \bibnamefont {Hou}},
  \bibinfo {author} {\bibfnamefont {J.}~\bibnamefont {Gao}}, \bibinfo {author}
  {\bibfnamefont {M.}~\bibnamefont {Guzzi}}, \bibinfo {author} {\bibfnamefont
  {J.}~\bibnamefont {Huston}}, \bibinfo {author} {\bibfnamefont
  {P.}~\bibnamefont {Nadolsky}}, \bibinfo {author} {\bibfnamefont
  {J.}~\bibnamefont {Pumplin}}, \bibinfo {author} {\bibfnamefont
  {C.}~\bibnamefont {Schmidt}}, \bibinfo {author} {\bibfnamefont
  {D.}~\bibnamefont {Stump}}, \ and\ \bibinfo {author} {\bibfnamefont {C.~P.}\
  \bibnamefont {Yuan}},\ }\href@noop {} {\bibfield  {journal} {\bibinfo
  {journal} {Phys. Rev.}\ }\textbf {\bibinfo {volume} {D93}},\ \bibinfo {pages}
  {033006} (\bibinfo {year} {2016})},\ \Eprint
  {http://arxiv.org/abs/1506.07443} {arXiv:1506.07443 [hep-ph]} \BibitemShut
  {NoStop}%
%%CITATION = ARXIV:1506.07443;%%
\bibitem [{\citenamefont {Martin}\ \emph {et~al.}(2009)\citenamefont {Martin},
  \citenamefont {Stirling}, \citenamefont {Thorne},\ and\ \citenamefont
  {Watt}}]{Martin:2009iq}%
  \BibitemOpen
  \bibfield  {author} {\bibinfo {author} {\bibfnamefont {A.~D.}\ \bibnamefont
  {Martin}}, \bibinfo {author} {\bibfnamefont {W.~J.}\ \bibnamefont
  {Stirling}}, \bibinfo {author} {\bibfnamefont {R.~S.}\ \bibnamefont
  {Thorne}}, \ and\ \bibinfo {author} {\bibfnamefont {G.}~\bibnamefont
  {Watt}},\ }\href@noop {} {\bibfield  {journal} {\bibinfo  {journal} {Eur.
  Phys. J.}\ }\textbf {\bibinfo {volume} {C63}},\ \bibinfo {pages} {189}
  (\bibinfo {year} {2009})},\ \Eprint {http://arxiv.org/abs/0901.0002}
  {arXiv:0901.0002 [hep-ph]} \BibitemShut {NoStop}%
%%CITATION = ARXIV:0901.0002;%%
\bibitem [{\citenamefont {Alekhin}\ \emph {et~al.}(2017)\citenamefont
  {Alekhin}, \citenamefont {Bl{\"u}mlein}, \citenamefont {Moch},\ and\
  \citenamefont {Placakyte}}]{Alekhin:2017kpj}%
  \BibitemOpen
  \bibfield  {author} {\bibinfo {author} {\bibfnamefont {S.}~\bibnamefont
  {Alekhin}}, \bibinfo {author} {\bibfnamefont {J.}~\bibnamefont
  {Bl{\"u}mlein}}, \bibinfo {author} {\bibfnamefont {S.}~\bibnamefont {Moch}},
  \ and\ \bibinfo {author} {\bibfnamefont {R.}~\bibnamefont {Placakyte}},\
  }\href@noop {} {\  (\bibinfo {year} {2017})},\ \Eprint
  {http://arxiv.org/abs/1701.05838} {arXiv:1701.05838 [hep-ph]} \BibitemShut
  {NoStop}%
%%CITATION = ARXIV:1701.05838;%%
\bibitem [{\citenamefont {Buckley}\ \emph {et~al.}(2015)\citenamefont
  {Buckley}, \citenamefont {Ferrando}, \citenamefont {Lloyd}, \citenamefont
  {Nordstr{\"o}m}, \citenamefont {Page}, \citenamefont {R{\"u}fenacht},
  \citenamefont {Sch{\"o}nherr},\ and\ \citenamefont {Watt}}]{Buckley:2014ana}%
  \BibitemOpen
  \bibfield  {author} {\bibinfo {author} {\bibfnamefont {A.}~\bibnamefont
  {Buckley}}, \bibinfo {author} {\bibfnamefont {J.}~\bibnamefont {Ferrando}},
  \bibinfo {author} {\bibfnamefont {S.}~\bibnamefont {Lloyd}}, \bibinfo
  {author} {\bibfnamefont {K.}~\bibnamefont {Nordstr{\"o}m}}, \bibinfo {author}
  {\bibfnamefont {B.}~\bibnamefont {Page}}, \bibinfo {author} {\bibfnamefont
  {M.}~\bibnamefont {R{\"u}fenacht}}, \bibinfo {author} {\bibfnamefont
  {M.}~\bibnamefont {Sch{\"o}nherr}}, \ and\ \bibinfo {author} {\bibfnamefont
  {G.}~\bibnamefont {Watt}},\ }\href@noop {} {\bibfield  {journal} {\bibinfo
  {journal} {Eur. Phys. J.}\ }\textbf {\bibinfo {volume} {C75}},\ \bibinfo
  {pages} {132} (\bibinfo {year} {2015})},\ \Eprint
  {http://arxiv.org/abs/1412.7420} {arXiv:1412.7420 [hep-ph]} \BibitemShut
  {NoStop}%
%%CITATION = ARXIV:1412.7420;%%
\bibitem [{\citenamefont {Martinelli}\ and\ \citenamefont
  {Sachrajda}(1987)}]{Martinelli:1987zd}%
  \BibitemOpen
  \bibfield  {author} {\bibinfo {author} {\bibfnamefont {G.}~\bibnamefont
  {Martinelli}}\ and\ \bibinfo {author} {\bibfnamefont {C.~T.}\ \bibnamefont
  {Sachrajda}},\ }\href@noop {} {\bibfield  {journal} {\bibinfo  {journal}
  {Phys. Lett.}\ }\textbf {\bibinfo {volume} {B196}},\ \bibinfo {pages} {184}
  (\bibinfo {year} {1987})}\BibitemShut {NoStop}%
%%CITATION = PHLTA,B196,184;%%
\bibitem [{\citenamefont {Martinelli}\ and\ \citenamefont
  {Sachrajda}(1989)}]{Martinelli:1988xs}%
  \BibitemOpen
  \bibfield  {author} {\bibinfo {author} {\bibfnamefont {G.}~\bibnamefont
  {Martinelli}}\ and\ \bibinfo {author} {\bibfnamefont {C.~T.}\ \bibnamefont
  {Sachrajda}},\ }\href@noop {} {\bibfield  {journal} {\bibinfo  {journal}
  {Phys. Lett.}\ }\textbf {\bibinfo {volume} {B217}},\ \bibinfo {pages} {319}
  (\bibinfo {year} {1989})}\BibitemShut {NoStop}%
%%CITATION = PHLTA,B217,319;%%
\bibitem [{\citenamefont {Detmold}\ \emph {et~al.}(2001)\citenamefont
  {Detmold}, \citenamefont {Melnitchouk},\ and\ \citenamefont
  {Thomas}}]{Detmold:2001dv}%
  \BibitemOpen
  \bibfield  {author} {\bibinfo {author} {\bibfnamefont {W.}~\bibnamefont
  {Detmold}}, \bibinfo {author} {\bibfnamefont {W.}~\bibnamefont
  {Melnitchouk}}, \ and\ \bibinfo {author} {\bibfnamefont {A.~W.}\ \bibnamefont
  {Thomas}},\ }\href@noop {} {\bibfield  {journal} {\bibinfo  {journal} {Eur.
  Phys. J.direct}\ }\textbf {\bibinfo {volume} {3}},\ \bibinfo {pages} {13}
  (\bibinfo {year} {2001})},\ \Eprint {http://arxiv.org/abs/hep-lat/0108002}
  {arXiv:hep-lat/0108002 [hep-lat]} \BibitemShut {NoStop}%
%%CITATION = HEP-LAT/0108002;%%
\bibitem [{\citenamefont {Detmold}\ \emph {et~al.}(2002)\citenamefont
  {Detmold}, \citenamefont {Melnitchouk},\ and\ \citenamefont
  {Thomas}}]{Detmold:2002nf}%
  \BibitemOpen
  \bibfield  {author} {\bibinfo {author} {\bibfnamefont {W.}~\bibnamefont
  {Detmold}}, \bibinfo {author} {\bibfnamefont {W.}~\bibnamefont
  {Melnitchouk}}, \ and\ \bibinfo {author} {\bibfnamefont {A.~W.}\ \bibnamefont
  {Thomas}},\ }\href@noop {} {\bibfield  {journal} {\bibinfo  {journal} {Phys.
  Rev.}\ }\textbf {\bibinfo {volume} {D66}},\ \bibinfo {pages} {054501}
  (\bibinfo {year} {2002})},\ \Eprint {http://arxiv.org/abs/hep-lat/0206001}
  {arXiv:hep-lat/0206001 [hep-lat]} \BibitemShut {NoStop}%
%%CITATION = HEP-LAT/0206001;%%
\bibitem [{\citenamefont {Dolgov}\ \emph {et~al.}(2002)\citenamefont {Dolgov}
  \emph {et~al.}}]{Dolgov:2002zm}%
  \BibitemOpen
  \bibfield  {author} {\bibinfo {author} {\bibfnamefont {D.}~\bibnamefont
  {Dolgov}} \emph {et~al.} (\bibinfo {collaboration} {LHPC, TXL}),\ }\href@noop
  {} {\bibfield  {journal} {\bibinfo  {journal} {Phys. Rev.}\ }\textbf
  {\bibinfo {volume} {D66}},\ \bibinfo {pages} {034506} (\bibinfo {year}
  {2002})},\ \Eprint {http://arxiv.org/abs/hep-lat/0201021}
  {arXiv:hep-lat/0201021 [hep-lat]} \BibitemShut {NoStop}%
%%CITATION = HEP-LAT/0201021;%%
\bibitem [{\citenamefont {Davoudi}\ and\ \citenamefont
  {Savage}(2012)}]{Davoudi:2012ya}%
  \BibitemOpen
  \bibfield  {author} {\bibinfo {author} {\bibfnamefont {Z.}~\bibnamefont
  {Davoudi}}\ and\ \bibinfo {author} {\bibfnamefont {M.~J.}\ \bibnamefont
  {Savage}},\ }\href@noop {} {\bibfield  {journal} {\bibinfo  {journal} {Phys.
  Rev.}\ }\textbf {\bibinfo {volume} {D86}},\ \bibinfo {pages} {054505}
  (\bibinfo {year} {2012})},\ \Eprint {http://arxiv.org/abs/1204.4146}
  {arXiv:1204.4146 [hep-lat]} \BibitemShut {NoStop}%
%%CITATION = ARXIV:1204.4146;%%
\bibitem [{\citenamefont {Liu}\ and\ \citenamefont {Dong}(1994)}]{Liu:1993cv}%
  \BibitemOpen
  \bibfield  {author} {\bibinfo {author} {\bibfnamefont {K.-F.}\ \bibnamefont
  {Liu}}\ and\ \bibinfo {author} {\bibfnamefont {S.-J.}\ \bibnamefont {Dong}},\
  }\href@noop {} {\bibfield  {journal} {\bibinfo  {journal} {Phys. Rev. Lett.}\
  }\textbf {\bibinfo {volume} {72}},\ \bibinfo {pages} {1790} (\bibinfo {year}
  {1994})},\ \Eprint {http://arxiv.org/abs/hep-ph/9306299}
  {arXiv:hep-ph/9306299 [hep-ph]} \BibitemShut {NoStop}%
%%CITATION = HEP-PH/9306299;%%
\bibitem [{\citenamefont {Liu}\ \emph {et~al.}(1999)\citenamefont {Liu},
  \citenamefont {Dong}, \citenamefont {Draper}, \citenamefont {Leinweber},
  \citenamefont {Sloan}, \citenamefont {Wilcox},\ and\ \citenamefont
  {Woloshyn}}]{Liu:1998um}%
  \BibitemOpen
  \bibfield  {author} {\bibinfo {author} {\bibfnamefont {K.~F.}\ \bibnamefont
  {Liu}}, \bibinfo {author} {\bibfnamefont {S.~J.}\ \bibnamefont {Dong}},
  \bibinfo {author} {\bibfnamefont {T.}~\bibnamefont {Draper}}, \bibinfo
  {author} {\bibfnamefont {D.}~\bibnamefont {Leinweber}}, \bibinfo {author}
  {\bibfnamefont {J.~H.}\ \bibnamefont {Sloan}}, \bibinfo {author}
  {\bibfnamefont {W.}~\bibnamefont {Wilcox}}, \ and\ \bibinfo {author}
  {\bibfnamefont {R.~M.}\ \bibnamefont {Woloshyn}},\ }\href@noop {} {\bibfield
  {journal} {\bibinfo  {journal} {Phys. Rev.}\ }\textbf {\bibinfo {volume}
  {D59}},\ \bibinfo {pages} {112001} (\bibinfo {year} {1999})},\ \Eprint
  {http://arxiv.org/abs/hep-ph/9806491} {arXiv:hep-ph/9806491 [hep-ph]}
  \BibitemShut {NoStop}%
%%CITATION = HEP-PH/9806491;%%
\bibitem [{\citenamefont {Liu}(2000)}]{Liu:1999ak}%
  \BibitemOpen
  \bibfield  {author} {\bibinfo {author} {\bibfnamefont {K.-F.}\ \bibnamefont
  {Liu}},\ }\href@noop {} {\bibfield  {journal} {\bibinfo  {journal} {Phys.
  Rev.}\ }\textbf {\bibinfo {volume} {D62}},\ \bibinfo {pages} {074501}
  (\bibinfo {year} {2000})},\ \Eprint {http://arxiv.org/abs/hep-ph/9910306}
  {arXiv:hep-ph/9910306 [hep-ph]} \BibitemShut {NoStop}%
%%CITATION = HEP-PH/9910306;%%
\bibitem [{\citenamefont {Liu}(2016)}]{Liu:2016djw}%
  \BibitemOpen
  \bibfield  {author} {\bibinfo {author} {\bibfnamefont {K.-F.}\ \bibnamefont
  {Liu}},\ }\bibfield  {booktitle} {\emph {\bibinfo {booktitle} {{Proceedings,
  33rd International Symposium on Lattice Field Theory (Lattice 2015): Kobe,
  Japan, July 14-18, 2015}}},\ }\href@noop {} {\bibfield  {journal} {\bibinfo
  {journal} {PoS}\ }\textbf {\bibinfo {volume} {LATTICE2015}},\ \bibinfo
  {pages} {115} (\bibinfo {year} {2016})},\ \Eprint
  {http://arxiv.org/abs/1603.07352} {arXiv:1603.07352 [hep-ph]} \BibitemShut
  {NoStop}%
%%CITATION = ARXIV:1603.07352;%%
\bibitem [{\citenamefont {Chambers}\ \emph {et~al.}(2017)\citenamefont
  {Chambers}, \citenamefont {Horsley}, \citenamefont {Nakamura}, \citenamefont
  {Perlt}, \citenamefont {Rakow}, \citenamefont {Schierholz}, \citenamefont
  {Schiller}, \citenamefont {Somfleth}, \citenamefont {Young},\ and\
  \citenamefont {Zanotti}}]{Chambers:2017dov}%
  \BibitemOpen
  \bibfield  {author} {\bibinfo {author} {\bibfnamefont {A.~J.}\ \bibnamefont
  {Chambers}}, \bibinfo {author} {\bibfnamefont {R.}~\bibnamefont {Horsley}},
  \bibinfo {author} {\bibfnamefont {Y.}~\bibnamefont {Nakamura}}, \bibinfo
  {author} {\bibfnamefont {H.}~\bibnamefont {Perlt}}, \bibinfo {author}
  {\bibfnamefont {P.~E.~L.}\ \bibnamefont {Rakow}}, \bibinfo {author}
  {\bibfnamefont {G.}~\bibnamefont {Schierholz}}, \bibinfo {author}
  {\bibfnamefont {A.}~\bibnamefont {Schiller}}, \bibinfo {author}
  {\bibfnamefont {K.}~\bibnamefont {Somfleth}}, \bibinfo {author}
  {\bibfnamefont {R.~D.}\ \bibnamefont {Young}}, \ and\ \bibinfo {author}
  {\bibfnamefont {J.~M.}\ \bibnamefont {Zanotti}},\ }\href@noop {} {\bibfield
  {journal} {\bibinfo  {journal} {Phys. Rev. Lett.}\ }\textbf {\bibinfo
  {volume} {118}},\ \bibinfo {pages} {242001} (\bibinfo {year} {2017})},\
  \Eprint {http://arxiv.org/abs/1703.01153} {arXiv:1703.01153 [hep-lat]}
  \BibitemShut {NoStop}%
%%CITATION = ARXIV:1703.01153;%%
\bibitem [{\citenamefont {Detmold}\ and\ \citenamefont
  {Lin}(2006)}]{Detmold:2005gg}%
  \BibitemOpen
  \bibfield  {author} {\bibinfo {author} {\bibfnamefont {W.}~\bibnamefont
  {Detmold}}\ and\ \bibinfo {author} {\bibfnamefont {C.~J.~D.}\ \bibnamefont
  {Lin}},\ }\href@noop {} {\bibfield  {journal} {\bibinfo  {journal} {Phys.
  Rev.}\ }\textbf {\bibinfo {volume} {D73}},\ \bibinfo {pages} {014501}
  (\bibinfo {year} {2006})},\ \Eprint {http://arxiv.org/abs/hep-lat/0507007}
  {arXiv:hep-lat/0507007 [hep-lat]} \BibitemShut {NoStop}%
%%CITATION = HEP-LAT/0507007;%%
\bibitem [{\citenamefont {Ma}\ and\ \citenamefont {Qiu}(2014)}]{Ma:2014jla}%
  \BibitemOpen
  \bibfield  {author} {\bibinfo {author} {\bibfnamefont {Y.-Q.}\ \bibnamefont
  {Ma}}\ and\ \bibinfo {author} {\bibfnamefont {J.-W.}\ \bibnamefont {Qiu}},\
  }\href@noop {} {\  (\bibinfo {year} {2014})},\ \Eprint
  {http://arxiv.org/abs/1404.6860} {arXiv:1404.6860 [hep-ph]} \BibitemShut
  {NoStop}%
%%CITATION = ARXIV:1404.6860;%%
\bibitem [{\citenamefont {Ma}\ and\ \citenamefont {Qiu}(2017)}]{Ma:2017pxb}%
  \BibitemOpen
  \bibfield  {author} {\bibinfo {author} {\bibfnamefont {Y.-Q.}\ \bibnamefont
  {Ma}}\ and\ \bibinfo {author} {\bibfnamefont {J.-W.}\ \bibnamefont {Qiu}},\
  }\href@noop {} {\  (\bibinfo {year} {2017})},\ \Eprint
  {http://arxiv.org/abs/1709.03018} {arXiv:1709.03018 [hep-ph]} \BibitemShut
  {NoStop}%
%%CITATION = ARXIV:1709.03018;%%
\bibitem [{\citenamefont {Ji}(2013)}]{Ji:2013dva}%
  \BibitemOpen
  \bibfield  {author} {\bibinfo {author} {\bibfnamefont {X.}~\bibnamefont
  {Ji}},\ }\href@noop {} {\bibfield  {journal} {\bibinfo  {journal} {Phys. Rev.
  Lett.}\ }\textbf {\bibinfo {volume} {110}},\ \bibinfo {pages} {262002}
  (\bibinfo {year} {2013})},\ \Eprint {http://arxiv.org/abs/1305.1539}
  {arXiv:1305.1539 [hep-ph]} \BibitemShut {NoStop}%
%%CITATION = ARXIV:1305.1539;%%
\bibitem [{\citenamefont {Ji}(2014)}]{Ji:2014gla}%
  \BibitemOpen
  \bibfield  {author} {\bibinfo {author} {\bibfnamefont {X.}~\bibnamefont
  {Ji}},\ }\href@noop {} {\bibfield  {journal} {\bibinfo  {journal} {Sci. China
  Phys. Mech. Astron.}\ }\textbf {\bibinfo {volume} {57}},\ \bibinfo {pages}
  {1407} (\bibinfo {year} {2014})},\ \Eprint {http://arxiv.org/abs/1404.6680}
  {arXiv:1404.6680 [hep-ph]} \BibitemShut {NoStop}%
%%CITATION = ARXIV:1404.6680;%%
\bibitem [{\citenamefont {Hatta}\ \emph {et~al.}(2014)\citenamefont {Hatta},
  \citenamefont {Ji},\ and\ \citenamefont {Zhao}}]{Hatta:2013gta}%
  \BibitemOpen
  \bibfield  {author} {\bibinfo {author} {\bibfnamefont {Y.}~\bibnamefont
  {Hatta}}, \bibinfo {author} {\bibfnamefont {X.}~\bibnamefont {Ji}}, \ and\
  \bibinfo {author} {\bibfnamefont {Y.}~\bibnamefont {Zhao}},\ }\href@noop {}
  {\bibfield  {journal} {\bibinfo  {journal} {Phys. Rev.}\ }\textbf {\bibinfo
  {volume} {D89}},\ \bibinfo {pages} {085030} (\bibinfo {year} {2014})},\
  \Eprint {http://arxiv.org/abs/1310.4263} {arXiv:1310.4263 [hep-ph]}
  \BibitemShut {NoStop}%
%%CITATION = ARXIV:1310.4263;%%
\bibitem [{\citenamefont {Xiong}\ \emph {et~al.}(2014)\citenamefont {Xiong},
  \citenamefont {Ji}, \citenamefont {Zhang},\ and\ \citenamefont
  {Zhao}}]{Xiong:2013bka}%
  \BibitemOpen
  \bibfield  {author} {\bibinfo {author} {\bibfnamefont {X.}~\bibnamefont
  {Xiong}}, \bibinfo {author} {\bibfnamefont {X.}~\bibnamefont {Ji}}, \bibinfo
  {author} {\bibfnamefont {J.-H.}\ \bibnamefont {Zhang}}, \ and\ \bibinfo
  {author} {\bibfnamefont {Y.}~\bibnamefont {Zhao}},\ }\href@noop {} {\bibfield
   {journal} {\bibinfo  {journal} {Phys. Rev.}\ }\textbf {\bibinfo {volume}
  {D90}},\ \bibinfo {pages} {014051} (\bibinfo {year} {2014})},\ \Eprint
  {http://arxiv.org/abs/1310.7471} {arXiv:1310.7471 [hep-ph]} \BibitemShut
  {NoStop}%
%%CITATION = ARXIV:1310.7471;%%
\bibitem [{\citenamefont {Alexandrou}\ \emph {et~al.}(2015)\citenamefont
  {Alexandrou}, \citenamefont {Cichy}, \citenamefont {Drach}, \citenamefont
  {Garcia-Ramos}, \citenamefont {Hadjiyiannakou}, \citenamefont {Jansen},
  \citenamefont {Steffens},\ and\ \citenamefont {Wiese}}]{Alexandrou:2015rja}%
  \BibitemOpen
  \bibfield  {author} {\bibinfo {author} {\bibfnamefont {C.}~\bibnamefont
  {Alexandrou}}, \bibinfo {author} {\bibfnamefont {K.}~\bibnamefont {Cichy}},
  \bibinfo {author} {\bibfnamefont {V.}~\bibnamefont {Drach}}, \bibinfo
  {author} {\bibfnamefont {E.}~\bibnamefont {Garcia-Ramos}}, \bibinfo {author}
  {\bibfnamefont {K.}~\bibnamefont {Hadjiyiannakou}}, \bibinfo {author}
  {\bibfnamefont {K.}~\bibnamefont {Jansen}}, \bibinfo {author} {\bibfnamefont
  {F.}~\bibnamefont {Steffens}}, \ and\ \bibinfo {author} {\bibfnamefont
  {C.}~\bibnamefont {Wiese}},\ }\href@noop {} {\bibfield  {journal} {\bibinfo
  {journal} {Phys. Rev.}\ }\textbf {\bibinfo {volume} {D92}},\ \bibinfo {pages}
  {014502} (\bibinfo {year} {2015})},\ \Eprint
  {http://arxiv.org/abs/1504.07455} {arXiv:1504.07455 [hep-lat]} \BibitemShut
  {NoStop}%
%%CITATION = ARXIV:1504.07455;%%
\bibitem [{\citenamefont {Stewart}\ and\ \citenamefont
  {Zhao}(2017)}]{Stewart:2017tvs}%
  \BibitemOpen
  \bibfield  {author} {\bibinfo {author} {\bibfnamefont {I.~W.}\ \bibnamefont
  {Stewart}}\ and\ \bibinfo {author} {\bibfnamefont {Y.}~\bibnamefont {Zhao}},\
  }\href@noop {} {\  (\bibinfo {year} {2017})},\ \Eprint
  {http://arxiv.org/abs/1709.04933} {arXiv:1709.04933 [hep-ph]} \BibitemShut
  {NoStop}%
%%CITATION = ARXIV:1709.04933;%%
\bibitem [{\citenamefont {Wang}\ \emph {et~al.}(2017)\citenamefont {Wang},
  \citenamefont {Zhao},\ and\ \citenamefont {Zhu}}]{Wang:2017qyg}%
  \BibitemOpen
  \bibfield  {author} {\bibinfo {author} {\bibfnamefont {W.}~\bibnamefont
  {Wang}}, \bibinfo {author} {\bibfnamefont {S.}~\bibnamefont {Zhao}}, \ and\
  \bibinfo {author} {\bibfnamefont {R.}~\bibnamefont {Zhu}},\ }\href@noop {} {\
   (\bibinfo {year} {2017})},\ \Eprint {http://arxiv.org/abs/1708.02458}
  {arXiv:1708.02458 [hep-ph]} \BibitemShut {NoStop}%
%%CITATION = ARXIV:1708.02458;%%
\bibitem [{\citenamefont {Wang}\ and\ \citenamefont
  {Zhao}(2017)}]{Wang:2017eel}%
  \BibitemOpen
  \bibfield  {author} {\bibinfo {author} {\bibfnamefont {W.}~\bibnamefont
  {Wang}}\ and\ \bibinfo {author} {\bibfnamefont {S.}~\bibnamefont {Zhao}},\
  }\href@noop {} {\  (\bibinfo {year} {2017})},\ \Eprint
  {http://arxiv.org/abs/1712.09247} {arXiv:1712.09247 [hep-ph]} \BibitemShut
  {NoStop}%
%%CITATION = ARXIV:1712.09247;%%
\bibitem [{\citenamefont {Lin}\ \emph {et~al.}(2015)\citenamefont {Lin},
  \citenamefont {Chen}, \citenamefont {Cohen},\ and\ \citenamefont
  {Ji}}]{Lin:2014zya}%
  \BibitemOpen
  \bibfield  {author} {\bibinfo {author} {\bibfnamefont {H.-W.}\ \bibnamefont
  {Lin}}, \bibinfo {author} {\bibfnamefont {J.-W.}\ \bibnamefont {Chen}},
  \bibinfo {author} {\bibfnamefont {S.~D.}\ \bibnamefont {Cohen}}, \ and\
  \bibinfo {author} {\bibfnamefont {X.}~\bibnamefont {Ji}},\ }\href@noop {}
  {\bibfield  {journal} {\bibinfo  {journal} {Phys. Rev.}\ }\textbf {\bibinfo
  {volume} {D91}},\ \bibinfo {pages} {054510} (\bibinfo {year} {2015})},\
  \Eprint {http://arxiv.org/abs/1402.1462} {arXiv:1402.1462 [hep-ph]}
  \BibitemShut {NoStop}%
%%CITATION = ARXIV:1402.1462;%%
\bibitem [{\citenamefont {Chen}\ \emph {et~al.}(2016)\citenamefont {Chen},
  \citenamefont {Cohen}, \citenamefont {Ji}, \citenamefont {Lin},\ and\
  \citenamefont {Zhang}}]{Chen:2016utp}%
  \BibitemOpen
  \bibfield  {author} {\bibinfo {author} {\bibfnamefont {J.-W.}\ \bibnamefont
  {Chen}}, \bibinfo {author} {\bibfnamefont {S.~D.}\ \bibnamefont {Cohen}},
  \bibinfo {author} {\bibfnamefont {X.}~\bibnamefont {Ji}}, \bibinfo {author}
  {\bibfnamefont {H.-W.}\ \bibnamefont {Lin}}, \ and\ \bibinfo {author}
  {\bibfnamefont {J.-H.}\ \bibnamefont {Zhang}},\ }\href@noop {} {\bibfield
  {journal} {\bibinfo  {journal} {Nucl. Phys.}\ }\textbf {\bibinfo {volume}
  {B911}},\ \bibinfo {pages} {246} (\bibinfo {year} {2016})},\ \Eprint
  {http://arxiv.org/abs/1603.06664} {arXiv:1603.06664 [hep-ph]} \BibitemShut
  {NoStop}%
%%CITATION = ARXIV:1603.06664;%%
\bibitem [{\citenamefont {Alexandrou}\ \emph
  {et~al.}(2017{\natexlab{a}})\citenamefont {Alexandrou}, \citenamefont
  {Cichy}, \citenamefont {Constantinou}, \citenamefont {Hadjiyiannakou},
  \citenamefont {Jansen}, \citenamefont {Steffens},\ and\ \citenamefont
  {Wiese}}]{Alexandrou:2016jqi}%
  \BibitemOpen
  \bibfield  {author} {\bibinfo {author} {\bibfnamefont {C.}~\bibnamefont
  {Alexandrou}}, \bibinfo {author} {\bibfnamefont {K.}~\bibnamefont {Cichy}},
  \bibinfo {author} {\bibfnamefont {M.}~\bibnamefont {Constantinou}}, \bibinfo
  {author} {\bibfnamefont {K.}~\bibnamefont {Hadjiyiannakou}}, \bibinfo
  {author} {\bibfnamefont {K.}~\bibnamefont {Jansen}}, \bibinfo {author}
  {\bibfnamefont {F.}~\bibnamefont {Steffens}}, \ and\ \bibinfo {author}
  {\bibfnamefont {C.}~\bibnamefont {Wiese}},\ }\href@noop {} {\bibfield
  {journal} {\bibinfo  {journal} {Phys. Rev.}\ }\textbf {\bibinfo {volume}
  {D96}},\ \bibinfo {pages} {014513} (\bibinfo {year} {2017}{\natexlab{a}})},\
  \Eprint {http://arxiv.org/abs/1610.03689} {arXiv:1610.03689 [hep-lat]}
  \BibitemShut {NoStop}%
%%CITATION = ARXIV:1610.03689;%%
\bibitem [{\citenamefont {Zhang}\ \emph {et~al.}(2017)\citenamefont {Zhang},
  \citenamefont {Chen}, \citenamefont {Ji}, \citenamefont {Jin},\ and\
  \citenamefont {Lin}}]{Zhang:2017bzy}%
  \BibitemOpen
  \bibfield  {author} {\bibinfo {author} {\bibfnamefont {J.-H.}\ \bibnamefont
  {Zhang}}, \bibinfo {author} {\bibfnamefont {J.-W.}\ \bibnamefont {Chen}},
  \bibinfo {author} {\bibfnamefont {X.}~\bibnamefont {Ji}}, \bibinfo {author}
  {\bibfnamefont {L.}~\bibnamefont {Jin}}, \ and\ \bibinfo {author}
  {\bibfnamefont {H.-W.}\ \bibnamefont {Lin}},\ }\href@noop {} {\bibfield
  {journal} {\bibinfo  {journal} {Phys. Rev.}\ }\textbf {\bibinfo {volume}
  {D95}},\ \bibinfo {pages} {094514} (\bibinfo {year} {2017})},\ \Eprint
  {http://arxiv.org/abs/1702.00008} {arXiv:1702.00008 [hep-lat]} \BibitemShut
  {NoStop}%
%%CITATION = ARXIV:1702.00008;%%
\bibitem [{\citenamefont {Alexandrou}\ \emph
  {et~al.}(2017{\natexlab{b}})\citenamefont {Alexandrou}, \citenamefont
  {Cichy}, \citenamefont {Constantinou}, \citenamefont {Hadjiyiannakou},
  \citenamefont {Jansen}, \citenamefont {Panagopoulos},\ and\ \citenamefont
  {Steffens}}]{Alexandrou:2017huk}%
  \BibitemOpen
  \bibfield  {author} {\bibinfo {author} {\bibfnamefont {C.}~\bibnamefont
  {Alexandrou}}, \bibinfo {author} {\bibfnamefont {K.}~\bibnamefont {Cichy}},
  \bibinfo {author} {\bibfnamefont {M.}~\bibnamefont {Constantinou}}, \bibinfo
  {author} {\bibfnamefont {K.}~\bibnamefont {Hadjiyiannakou}}, \bibinfo
  {author} {\bibfnamefont {K.}~\bibnamefont {Jansen}}, \bibinfo {author}
  {\bibfnamefont {H.}~\bibnamefont {Panagopoulos}}, \ and\ \bibinfo {author}
  {\bibfnamefont {F.}~\bibnamefont {Steffens}},\ }\href@noop {} {\bibfield
  {journal} {\bibinfo  {journal} {Nucl. Phys.}\ }\textbf {\bibinfo {volume}
  {B923}},\ \bibinfo {pages} {394} (\bibinfo {year} {2017}{\natexlab{b}})},\
  \Eprint {http://arxiv.org/abs/1706.00265} {arXiv:1706.00265 [hep-lat]}
  \BibitemShut {NoStop}%
%%CITATION = ARXIV:1706.00265;%%
\bibitem [{\citenamefont {Chen}\ \emph
  {et~al.}(2017{\natexlab{a}})\citenamefont {Chen}, \citenamefont {Ishikawa},
  \citenamefont {Jin}, \citenamefont {Lin}, \citenamefont {Yang}, \citenamefont
  {Zhang},\ and\ \citenamefont {Zhao}}]{Chen:2017mzz}%
  \BibitemOpen
  \bibfield  {author} {\bibinfo {author} {\bibfnamefont {J.-W.}\ \bibnamefont
  {Chen}}, \bibinfo {author} {\bibfnamefont {T.}~\bibnamefont {Ishikawa}},
  \bibinfo {author} {\bibfnamefont {L.}~\bibnamefont {Jin}}, \bibinfo {author}
  {\bibfnamefont {H.-W.}\ \bibnamefont {Lin}}, \bibinfo {author} {\bibfnamefont
  {Y.-B.}\ \bibnamefont {Yang}}, \bibinfo {author} {\bibfnamefont {J.-H.}\
  \bibnamefont {Zhang}}, \ and\ \bibinfo {author} {\bibfnamefont
  {Y.}~\bibnamefont {Zhao}},\ }\href@noop {} {\  (\bibinfo {year}
  {2017}{\natexlab{a}})},\ \Eprint {http://arxiv.org/abs/1706.01295}
  {arXiv:1706.01295 [hep-lat]} \BibitemShut {NoStop}%
%%CITATION = ARXIV:1706.01295;%%
\bibitem [{\citenamefont {Lin}\ \emph {et~al.}(2017{\natexlab{a}})\citenamefont
  {Lin}, \citenamefont {Chen}, \citenamefont {Ishikawa},\ and\ \citenamefont
  {Zhang}}]{Lin:2017ani}%
  \BibitemOpen
  \bibfield  {author} {\bibinfo {author} {\bibfnamefont {H.-W.}\ \bibnamefont
  {Lin}}, \bibinfo {author} {\bibfnamefont {J.-W.}\ \bibnamefont {Chen}},
  \bibinfo {author} {\bibfnamefont {T.}~\bibnamefont {Ishikawa}}, \ and\
  \bibinfo {author} {\bibfnamefont {J.-H.}\ \bibnamefont {Zhang}},\ }\href@noop
  {} {\  (\bibinfo {year} {2017}{\natexlab{a}})},\ \Eprint
  {http://arxiv.org/abs/1708.05301} {arXiv:1708.05301 [hep-lat]} \BibitemShut
  {NoStop}%
%%CITATION = ARXIV:1708.05301;%%
\bibitem [{\citenamefont {Chen}\ \emph
  {et~al.}(2017{\natexlab{b}})\citenamefont {Chen}, \citenamefont {Jin},
  \citenamefont {Lin}, \citenamefont {Sch{\"a}fer}, \citenamefont {Sun},
  \citenamefont {Yang}, \citenamefont {Zhang}, \citenamefont {Zhang},\ and\
  \citenamefont {Zhao}}]{Chen:2017gck}%
  \BibitemOpen
  \bibfield  {author} {\bibinfo {author} {\bibfnamefont {J.-W.}\ \bibnamefont
  {Chen}}, \bibinfo {author} {\bibfnamefont {L.}~\bibnamefont {Jin}}, \bibinfo
  {author} {\bibfnamefont {H.-W.}\ \bibnamefont {Lin}}, \bibinfo {author}
  {\bibfnamefont {A.}~\bibnamefont {Sch{\"a}fer}}, \bibinfo {author}
  {\bibfnamefont {P.}~\bibnamefont {Sun}}, \bibinfo {author} {\bibfnamefont
  {Y.-B.}\ \bibnamefont {Yang}}, \bibinfo {author} {\bibfnamefont {J.-H.}\
  \bibnamefont {Zhang}}, \bibinfo {author} {\bibfnamefont {R.}~\bibnamefont
  {Zhang}}, \ and\ \bibinfo {author} {\bibfnamefont {Y.}~\bibnamefont {Zhao}},\
  }\href@noop {} {\  (\bibinfo {year} {2017}{\natexlab{b}})},\ \Eprint
  {http://arxiv.org/abs/1712.10025} {arXiv:1712.10025 [hep-ph]} \BibitemShut
  {NoStop}%
%%CITATION = ARXIV:1712.10025;%%
\bibitem [{\citenamefont {Lin}\ \emph {et~al.}(2017{\natexlab{b}})\citenamefont
  {Lin} \emph {et~al.}}]{Lin:2017snn}%
  \BibitemOpen
  \bibfield  {author} {\bibinfo {author} {\bibfnamefont {H.-W.}\ \bibnamefont
  {Lin}} \emph {et~al.},\ }\href@noop {} {\  (\bibinfo {year}
  {2017}{\natexlab{b}})},\ \Eprint {http://arxiv.org/abs/1711.07916}
  {arXiv:1711.07916 [hep-ph]} \BibitemShut {NoStop}%
%%CITATION = ARXIV:1711.07916;%%
\bibitem [{\citenamefont {Ji}\ \emph {et~al.}(2015)\citenamefont {Ji},
  \citenamefont {Sch{\"a}fer}, \citenamefont {Xiong},\ and\ \citenamefont
  {Zhang}}]{Ji:2015qla}%
  \BibitemOpen
  \bibfield  {author} {\bibinfo {author} {\bibfnamefont {X.}~\bibnamefont
  {Ji}}, \bibinfo {author} {\bibfnamefont {A.}~\bibnamefont {Sch{\"a}fer}},
  \bibinfo {author} {\bibfnamefont {X.}~\bibnamefont {Xiong}}, \ and\ \bibinfo
  {author} {\bibfnamefont {J.-H.}\ \bibnamefont {Zhang}},\ }\href@noop {}
  {\bibfield  {journal} {\bibinfo  {journal} {Phys. Rev.}\ }\textbf {\bibinfo
  {volume} {D92}},\ \bibinfo {pages} {014039} (\bibinfo {year} {2015})},\
  \Eprint {http://arxiv.org/abs/1506.00248} {arXiv:1506.00248 [hep-ph]}
  \BibitemShut {NoStop}%
%%CITATION = ARXIV:1506.00248;%%
\bibitem [{\citenamefont
  {Radyushkin}(2017{\natexlab{a}})}]{Radyushkin:2017cyf}%
  \BibitemOpen
  \bibfield  {author} {\bibinfo {author} {\bibfnamefont {A.~V.}\ \bibnamefont
  {Radyushkin}},\ }\href@noop {} {\bibfield  {journal} {\bibinfo  {journal}
  {Phys. Rev.}\ }\textbf {\bibinfo {volume} {D96}},\ \bibinfo {pages} {034025}
  (\bibinfo {year} {2017}{\natexlab{a}})},\ \Eprint
  {http://arxiv.org/abs/1705.01488} {arXiv:1705.01488 [hep-ph]} \BibitemShut
  {NoStop}%
%%CITATION = ARXIV:1705.01488;%%
\bibitem [{\citenamefont {Radyushkin}(1983)}]{Radyushkin:1983wh}%
  \BibitemOpen
  \bibfield  {author} {\bibinfo {author} {\bibfnamefont {A.~V.}\ \bibnamefont
  {Radyushkin}},\ }\href@noop {} {\bibfield  {journal} {\bibinfo  {journal}
  {Phys. Lett.}\ }\textbf {\bibinfo {volume} {131B}},\ \bibinfo {pages} {179}
  (\bibinfo {year} {1983})}\BibitemShut {NoStop}%
%%CITATION = PHLTA,131B,179;%%
\bibitem [{\citenamefont
  {Radyushkin}(2017{\natexlab{b}})}]{Radyushkin:2016hsy}%
  \BibitemOpen
  \bibfield  {author} {\bibinfo {author} {\bibfnamefont {A.}~\bibnamefont
  {Radyushkin}},\ }\href@noop {} {\bibfield  {journal} {\bibinfo  {journal}
  {Phys. Lett.}\ }\textbf {\bibinfo {volume} {B767}},\ \bibinfo {pages} {314}
  (\bibinfo {year} {2017}{\natexlab{b}})},\ \Eprint
  {http://arxiv.org/abs/1612.05170} {arXiv:1612.05170 [hep-ph]} \BibitemShut
  {NoStop}%
%%CITATION = ARXIV:1612.05170;%%
\bibitem [{\citenamefont {Orginos}\ \emph {et~al.}(2017)\citenamefont
  {Orginos}, \citenamefont {Radyushkin}, \citenamefont {Karpie},\ and\
  \citenamefont {Zafeiropoulos}}]{Orginos:2017kos}%
  \BibitemOpen
  \bibfield  {author} {\bibinfo {author} {\bibfnamefont {K.}~\bibnamefont
  {Orginos}}, \bibinfo {author} {\bibfnamefont {A.}~\bibnamefont {Radyushkin}},
  \bibinfo {author} {\bibfnamefont {J.}~\bibnamefont {Karpie}}, \ and\ \bibinfo
  {author} {\bibfnamefont {S.}~\bibnamefont {Zafeiropoulos}},\ }\href@noop {}
  {\  (\bibinfo {year} {2017})},\ \Eprint {http://arxiv.org/abs/1706.05373}
  {arXiv:1706.05373 [hep-ph]} \BibitemShut {NoStop}%
%%CITATION = ARXIV:1706.05373;%%
\bibitem [{\citenamefont {Karpie}\ \emph {et~al.}(2017)\citenamefont {Karpie},
  \citenamefont {Orginos}, \citenamefont {Radyushkin},\ and\ \citenamefont
  {Zafeiropoulos}}]{Karpie:2017bzm}%
  \BibitemOpen
  \bibfield  {author} {\bibinfo {author} {\bibfnamefont {J.}~\bibnamefont
  {Karpie}}, \bibinfo {author} {\bibfnamefont {K.}~\bibnamefont {Orginos}},
  \bibinfo {author} {\bibfnamefont {A.}~\bibnamefont {Radyushkin}}, \ and\
  \bibinfo {author} {\bibfnamefont {S.}~\bibnamefont {Zafeiropoulos}},\
  }\href@noop {} {\  (\bibinfo {year} {2017})},\ \Eprint
  {http://arxiv.org/abs/1710.08288} {arXiv:1710.08288 [hep-lat]} \BibitemShut
  {NoStop}%
%%CITATION = ARXIV:1710.08288;%%
\bibitem [{\citenamefont {Ioffe}(1969)}]{Ioffe:1969kf}%
  \BibitemOpen
  \bibfield  {author} {\bibinfo {author} {\bibfnamefont {B.~L.}\ \bibnamefont
  {Ioffe}},\ }\href {\doibase 10.1016/0370-2693(69)90415-8} {\bibfield
  {journal} {\bibinfo  {journal} {Phys. Lett.}\ }\textbf {\bibinfo {volume}
  {30B}},\ \bibinfo {pages} {123} (\bibinfo {year} {1969})}\BibitemShut
  {NoStop}%
%%CITATION = PHLTA,30B,123;%%
\bibitem [{\citenamefont {Ji}\ \emph {et~al.}(2017{\natexlab{a}})\citenamefont
  {Ji}, \citenamefont {Zhang},\ and\ \citenamefont {Zhao}}]{Ji:2017rah}%
  \BibitemOpen
  \bibfield  {author} {\bibinfo {author} {\bibfnamefont {X.}~\bibnamefont
  {Ji}}, \bibinfo {author} {\bibfnamefont {J.-H.}\ \bibnamefont {Zhang}}, \
  and\ \bibinfo {author} {\bibfnamefont {Y.}~\bibnamefont {Zhao}},\ }\href@noop
  {} {\  (\bibinfo {year} {2017}{\natexlab{a}})},\ \Eprint
  {http://arxiv.org/abs/1706.07416} {arXiv:1706.07416 [hep-ph]} \BibitemShut
  {NoStop}%
%%CITATION = ARXIV:1706.07416;%%
\bibitem [{\citenamefont {Wilson}(1969)}]{Wilson:1969zs}%
  \BibitemOpen
  \bibfield  {author} {\bibinfo {author} {\bibfnamefont {K.~G.}\ \bibnamefont
  {Wilson}},\ }\href {\doibase 10.1103/PhysRev.179.1499} {\bibfield  {journal}
  {\bibinfo  {journal} {Phys. Rev.}\ }\textbf {\bibinfo {volume} {179}},\
  \bibinfo {pages} {1499} (\bibinfo {year} {1969})}\BibitemShut {NoStop}%
%%CITATION = PHRVA,179,1499;%%
\bibitem [{\citenamefont {Brandt}(1967)}]{Brandt:1967rb}%
  \BibitemOpen
  \bibfield  {author} {\bibinfo {author} {\bibfnamefont {R.~A.}\ \bibnamefont
  {Brandt}},\ }\href {\doibase 10.1016/0003-4916(67)90177-7} {\bibfield
  {journal} {\bibinfo  {journal} {Annals Phys.}\ }\textbf {\bibinfo {volume}
  {44}},\ \bibinfo {pages} {221} (\bibinfo {year} {1967})}\BibitemShut
  {NoStop}%
%%CITATION = APNYA,44,221;%%
\bibitem [{\citenamefont {Wilson}\ and\ \citenamefont
  {Zimmermann}(1972)}]{Wilson:1972ee}%
  \BibitemOpen
  \bibfield  {author} {\bibinfo {author} {\bibfnamefont {K.~G.}\ \bibnamefont
  {Wilson}}\ and\ \bibinfo {author} {\bibfnamefont {W.}~\bibnamefont
  {Zimmermann}},\ }\href@noop {} {\bibfield  {journal} {\bibinfo  {journal}
  {Commun. Math. Phys.}\ }\textbf {\bibinfo {volume} {24}},\ \bibinfo {pages}
  {87} (\bibinfo {year} {1972})}\BibitemShut {NoStop}%
%%CITATION = CMPHA,24,87;%%
\bibitem [{\citenamefont {Zimmermann}(1973)}]{Zimmermann:1972tv}%
  \BibitemOpen
  \bibfield  {author} {\bibinfo {author} {\bibfnamefont {W.}~\bibnamefont
  {Zimmermann}},\ }\bibfield  {booktitle} {\emph {\bibinfo {booktitle} {{In
  *Tegernsee 1998, Quantum field theory* 278-309}}},\ }\href@noop {} {\bibfield
   {journal} {\bibinfo  {journal} {Annals Phys.}\ }\textbf {\bibinfo {volume}
  {77}},\ \bibinfo {pages} {570} (\bibinfo {year} {1973})},\ \bibinfo {note}
  {[Lect. Notes Phys.558,278(2000)]}\BibitemShut {NoStop}%
%%CITATION = APNYA,77,570;%%
\bibitem [{\citenamefont {Dorn}(1986)}]{Dorn:1986dt}%
  \BibitemOpen
  \bibfield  {author} {\bibinfo {author} {\bibfnamefont {H.}~\bibnamefont
  {Dorn}},\ }\href@noop {} {\bibfield  {journal} {\bibinfo  {journal} {Fortsch.
  Phys.}\ }\textbf {\bibinfo {volume} {34}},\ \bibinfo {pages} {11} (\bibinfo
  {year} {1986})}\BibitemShut {NoStop}%
%%CITATION = FPYKA,34,11;%%
\bibitem [{\citenamefont {Zhao}(2017)}]{Zhao:2017int}%
  \BibitemOpen
  \bibfield  {author} {\bibinfo {author} {\bibfnamefont {Y.}~\bibnamefont
  {Zhao}},\ }\href
  {http://www.int.washington.edu/talks/WorkShops/int_17_68W/People/Zhao_Y/Zhao.pdf}
  {\emph {\bibinfo {title} {Factorization Theorem Relating Euclidean and
  Light-Cone Parton Distributions}}}\ (\bibinfo  {publisher} {Talk given at INT
  Workshop 17-68W, ``Flavor Structure of the Nucleon Sea"},\ \bibinfo {address}
  {Seattle, US},\ \bibinfo {year} {2017})\BibitemShut {NoStop}%
\bibitem [{\citenamefont {Ji}\ \emph {et~al.}(2017{\natexlab{b}})\citenamefont
  {Ji}, \citenamefont {Zhang},\ and\ \citenamefont {Zhao}}]{Ji:2017oey}%
  \BibitemOpen
  \bibfield  {author} {\bibinfo {author} {\bibfnamefont {X.}~\bibnamefont
  {Ji}}, \bibinfo {author} {\bibfnamefont {J.-H.}\ \bibnamefont {Zhang}}, \
  and\ \bibinfo {author} {\bibfnamefont {Y.}~\bibnamefont {Zhao}},\ }\href@noop
  {} {\  (\bibinfo {year} {2017}{\natexlab{b}})},\ \Eprint
  {http://arxiv.org/abs/1706.08962} {arXiv:1706.08962 [hep-ph]} \BibitemShut
  {NoStop}%
%%CITATION = ARXIV:1706.08962;%%
\bibitem [{\citenamefont {Ishikawa}\ \emph {et~al.}(2017)\citenamefont
  {Ishikawa}, \citenamefont {Ma}, \citenamefont {Qiu},\ and\ \citenamefont
  {Yoshida}}]{Ishikawa:2017faj}%
  \BibitemOpen
  \bibfield  {author} {\bibinfo {author} {\bibfnamefont {T.}~\bibnamefont
  {Ishikawa}}, \bibinfo {author} {\bibfnamefont {Y.-Q.}\ \bibnamefont {Ma}},
  \bibinfo {author} {\bibfnamefont {J.-W.}\ \bibnamefont {Qiu}}, \ and\
  \bibinfo {author} {\bibfnamefont {S.}~\bibnamefont {Yoshida}},\ }\href@noop
  {} {\  (\bibinfo {year} {2017})},\ \Eprint {http://arxiv.org/abs/1707.03107}
  {arXiv:1707.03107 [hep-ph]} \BibitemShut {NoStop}%
%%CITATION = ARXIV:1707.03107;%%
\bibitem [{\citenamefont {Ishikawa}\ \emph {et~al.}(2016)\citenamefont
  {Ishikawa}, \citenamefont {Ma}, \citenamefont {Qiu},\ and\ \citenamefont
  {Yoshida}}]{Ishikawa:2016znu}%
  \BibitemOpen
  \bibfield  {author} {\bibinfo {author} {\bibfnamefont {T.}~\bibnamefont
  {Ishikawa}}, \bibinfo {author} {\bibfnamefont {Y.-Q.}\ \bibnamefont {Ma}},
  \bibinfo {author} {\bibfnamefont {J.-W.}\ \bibnamefont {Qiu}}, \ and\
  \bibinfo {author} {\bibfnamefont {S.}~\bibnamefont {Yoshida}},\ }\href@noop
  {} {\  (\bibinfo {year} {2016})},\ \Eprint {http://arxiv.org/abs/1609.02018}
  {arXiv:1609.02018 [hep-lat]} \BibitemShut {NoStop}%
%%CITATION = ARXIV:1609.02018;%%
\bibitem [{\citenamefont {Chen}\ \emph
  {et~al.}(2017{\natexlab{c}})\citenamefont {Chen}, \citenamefont {Ji},\ and\
  \citenamefont {Zhang}}]{Chen:2016fxx}%
  \BibitemOpen
  \bibfield  {author} {\bibinfo {author} {\bibfnamefont {J.-W.}\ \bibnamefont
  {Chen}}, \bibinfo {author} {\bibfnamefont {X.}~\bibnamefont {Ji}}, \ and\
  \bibinfo {author} {\bibfnamefont {J.-H.}\ \bibnamefont {Zhang}},\ }\href@noop
  {} {\bibfield  {journal} {\bibinfo  {journal} {Nucl. Phys.}\ }\textbf
  {\bibinfo {volume} {B915}},\ \bibinfo {pages} {1} (\bibinfo {year}
  {2017}{\natexlab{c}})},\ \Eprint {http://arxiv.org/abs/1609.08102}
  {arXiv:1609.08102 [hep-ph]} \BibitemShut {NoStop}%
%%CITATION = ARXIV:1609.08102;%%
\bibitem [{\citenamefont {Constantinou}\ and\ \citenamefont
  {Panagopoulos}(2017)}]{Constantinou:2017sej}%
  \BibitemOpen
  \bibfield  {author} {\bibinfo {author} {\bibfnamefont {M.}~\bibnamefont
  {Constantinou}}\ and\ \bibinfo {author} {\bibfnamefont {H.}~\bibnamefont
  {Panagopoulos}},\ }\href@noop {} {\  (\bibinfo {year} {2017})},\ \Eprint
  {http://arxiv.org/abs/1705.11193} {arXiv:1705.11193 [hep-lat]} \BibitemShut
  {NoStop}%
%%CITATION = ARXIV:1705.11193;%%
\bibitem [{\citenamefont {Nachtmann}(1973)}]{Nachtmann:1973mr}%
  \BibitemOpen
  \bibfield  {author} {\bibinfo {author} {\bibfnamefont {O.}~\bibnamefont
  {Nachtmann}},\ }\href@noop {} {\bibfield  {journal} {\bibinfo  {journal}
  {Nucl. Phys.}\ }\textbf {\bibinfo {volume} {B63}},\ \bibinfo {pages} {237}
  (\bibinfo {year} {1973})}\BibitemShut {NoStop}%
%%CITATION = NUPHA,B63,237;%%
\bibitem [{\citenamefont {Georgi}\ and\ \citenamefont
  {Politzer}(1976)}]{Georgi:1976ve}%
  \BibitemOpen
  \bibfield  {author} {\bibinfo {author} {\bibfnamefont {H.}~\bibnamefont
  {Georgi}}\ and\ \bibinfo {author} {\bibfnamefont {H.~D.}\ \bibnamefont
  {Politzer}},\ }\href@noop {} {\bibfield  {journal} {\bibinfo  {journal}
  {Phys. Rev.}\ }\textbf {\bibinfo {volume} {D14}},\ \bibinfo {pages} {1829}
  (\bibinfo {year} {1976})}\BibitemShut {NoStop}%
%%CITATION = PHRVA,D14,1829;%%
\bibitem [{\citenamefont {Braun}\ and\ \citenamefont
  {Mueller}(2008)}]{Braun:2007wv}%
  \BibitemOpen
  \bibfield  {author} {\bibinfo {author} {\bibfnamefont {V.}~\bibnamefont
  {Braun}}\ and\ \bibinfo {author} {\bibfnamefont {D.}~\bibnamefont
  {Mueller}},\ }\href@noop {} {\bibfield  {journal} {\bibinfo  {journal} {Eur.
  Phys. J.}\ }\textbf {\bibinfo {volume} {C55}},\ \bibinfo {pages} {349}
  (\bibinfo {year} {2008})},\ \Eprint {http://arxiv.org/abs/0709.1348}
  {arXiv:0709.1348 [hep-ph]} \BibitemShut {NoStop}%
%%CITATION = ARXIV:0709.1348;%%
\bibitem [{\citenamefont {Bali}\ \emph {et~al.}(2017)\citenamefont {Bali},
  \citenamefont {Braun}, \citenamefont {G{\"o}ckeler}, \citenamefont {Gruber},
  \citenamefont {Hutzler}, \citenamefont {Korcyl}, \citenamefont {Lang},
  \citenamefont {Sch{\"a}fer}, \citenamefont {Wein},\ and\ \citenamefont
  {Zhang}}]{Bali:2017gfr}%
  \BibitemOpen
  \bibfield  {author} {\bibinfo {author} {\bibfnamefont {G.~S.}\ \bibnamefont
  {Bali}}, \bibinfo {author} {\bibfnamefont {V.~M.}\ \bibnamefont {Braun}},
  \bibinfo {author} {\bibfnamefont {M.}~\bibnamefont {G{\"o}ckeler}}, \bibinfo
  {author} {\bibfnamefont {M.}~\bibnamefont {Gruber}}, \bibinfo {author}
  {\bibfnamefont {F.}~\bibnamefont {Hutzler}}, \bibinfo {author} {\bibfnamefont
  {P.}~\bibnamefont {Korcyl}}, \bibinfo {author} {\bibfnamefont
  {B.}~\bibnamefont {Lang}}, \bibinfo {author} {\bibfnamefont {A.}~\bibnamefont
  {Sch{\"a}fer}}, \bibinfo {author} {\bibfnamefont {P.}~\bibnamefont {Wein}}, \
  and\ \bibinfo {author} {\bibfnamefont {J.-H.}\ \bibnamefont {Zhang}},\ }in\
  \href@noop {} {\emph {\bibinfo {booktitle} {{35th International Symposium on
  Lattice Field Theory (Lattice 2017) Granada, Spain, June 18-24, 2017}}}}\
  (\bibinfo {year} {2017})\ \Eprint {http://arxiv.org/abs/1709.04325}
  {arXiv:1709.04325 [hep-lat]} \BibitemShut {NoStop}%
%%CITATION = ARXIV:1709.04325;%%
\bibitem [{\citenamefont {Zhang}\ \emph {et~al.}(2018)\citenamefont {Zhang},
  \citenamefont {Chen},\ and\ \citenamefont {Monahan}}]{Zhang:2018ggy}%
  \BibitemOpen
  \bibfield  {author} {\bibinfo {author} {\bibfnamefont {J.-H.}\ \bibnamefont
  {Zhang}}, \bibinfo {author} {\bibfnamefont {J.-W.}\ \bibnamefont {Chen}}, \
  and\ \bibinfo {author} {\bibfnamefont {C.}~\bibnamefont {Monahan}},\
  }\href@noop {} {\  (\bibinfo {year} {2018})},\ \Eprint
  {http://arxiv.org/abs/1801.03023} {arXiv:1801.03023 [hep-ph]} \BibitemShut
  {NoStop}%
%%CITATION = ARXIV:1801.03023;%%
\bibitem [{\citenamefont
  {Radyushkin}(2017{\natexlab{c}})}]{Radyushkin:2017lvu}%
  \BibitemOpen
  \bibfield  {author} {\bibinfo {author} {\bibfnamefont {A.~V.}\ \bibnamefont
  {Radyushkin}},\ }\href@noop {} {\  (\bibinfo {year} {2017}{\natexlab{c}})},\
  \Eprint {http://arxiv.org/abs/1710.08813} {arXiv:1710.08813 [hep-ph]}
  \BibitemShut {NoStop}%
%%CITATION = ARXIV:1710.08813;%%
\bibitem [{\citenamefont {Radyushkin}(2018)}]{Radyushkin:2018cvn}%
  \BibitemOpen
  \bibfield  {author} {\bibinfo {author} {\bibfnamefont {A.}~\bibnamefont
  {Radyushkin}},\ }\href@noop {} {\  (\bibinfo {year} {2018})},\ \Eprint
  {http://arxiv.org/abs/1801.02427} {arXiv:1801.02427 [hep-ph]} \BibitemShut
  {NoStop}%
%%CITATION = ARXIV:1801.02427;%%
\bibitem [{\citenamefont {Ji}\ and\ \citenamefont {Zhang}(2015)}]{Ji:2015jwa}%
  \BibitemOpen
  \bibfield  {author} {\bibinfo {author} {\bibfnamefont {X.}~\bibnamefont
  {Ji}}\ and\ \bibinfo {author} {\bibfnamefont {J.-H.}\ \bibnamefont {Zhang}},\
  }\href@noop {} {\bibfield  {journal} {\bibinfo  {journal} {Phys. Rev.}\
  }\textbf {\bibinfo {volume} {D92}},\ \bibinfo {pages} {034006} (\bibinfo
  {year} {2015})},\ \Eprint {http://arxiv.org/abs/1505.07699} {arXiv:1505.07699
  [hep-ph]} \BibitemShut {NoStop}%
%%CITATION = ARXIV:1505.07699;%%
\bibitem [{\citenamefont {Xiong}\ \emph {et~al.}(2017)\citenamefont {Xiong},
  \citenamefont {Luu},\ and\ \citenamefont {Mei{\ss}ner}}]{Xiong:2017jtn}%
  \BibitemOpen
  \bibfield  {author} {\bibinfo {author} {\bibfnamefont {X.}~\bibnamefont
  {Xiong}}, \bibinfo {author} {\bibfnamefont {T.}~\bibnamefont {Luu}}, \ and\
  \bibinfo {author} {\bibfnamefont {U.-G.}\ \bibnamefont {Mei{\ss}ner}},\
  }\href@noop {} {\  (\bibinfo {year} {2017})},\ \Eprint
  {http://arxiv.org/abs/1705.00246} {arXiv:1705.00246 [hep-ph]} \BibitemShut
  {NoStop}%
%%CITATION = ARXIV:1705.00246;%%
\bibitem [{\citenamefont {Green}\ \emph {et~al.}(2017)\citenamefont {Green},
  \citenamefont {Jansen},\ and\ \citenamefont {Steffens}}]{Green:2017xeu}%
  \BibitemOpen
  \bibfield  {author} {\bibinfo {author} {\bibfnamefont {J.}~\bibnamefont
  {Green}}, \bibinfo {author} {\bibfnamefont {K.}~\bibnamefont {Jansen}}, \
  and\ \bibinfo {author} {\bibfnamefont {F.}~\bibnamefont {Steffens}},\
  }\href@noop {} {\  (\bibinfo {year} {2017})},\ \Eprint
  {http://arxiv.org/abs/1707.07152} {arXiv:1707.07152 [hep-lat]} \BibitemShut
  {NoStop}%
%%CITATION = ARXIV:1707.07152;%%
\bibitem [{\citenamefont {Chen}\ \emph
  {et~al.}(2017{\natexlab{d}})\citenamefont {Chen}, \citenamefont {Ishikawa},
  \citenamefont {Jin}, \citenamefont {Lin}, \citenamefont {Yang}, \citenamefont
  {Zhang},\ and\ \citenamefont {Zhao}}]{Chen:2017mie}%
  \BibitemOpen
  \bibfield  {author} {\bibinfo {author} {\bibfnamefont {J.-W.}\ \bibnamefont
  {Chen}}, \bibinfo {author} {\bibfnamefont {T.}~\bibnamefont {Ishikawa}},
  \bibinfo {author} {\bibfnamefont {L.}~\bibnamefont {Jin}}, \bibinfo {author}
  {\bibfnamefont {H.-W.}\ \bibnamefont {Lin}}, \bibinfo {author} {\bibfnamefont
  {Y.-B.}\ \bibnamefont {Yang}}, \bibinfo {author} {\bibfnamefont {J.-H.}\
  \bibnamefont {Zhang}}, \ and\ \bibinfo {author} {\bibfnamefont
  {Y.}~\bibnamefont {Zhao}},\ }\href@noop {} {\  (\bibinfo {year}
  {2017}{\natexlab{d}})},\ \Eprint {http://arxiv.org/abs/1710.01089}
  {arXiv:1710.01089 [hep-lat]} \BibitemShut {NoStop}%
%%CITATION = ARXIV:1710.01089;%%
\bibitem [{\citenamefont {Zhang}\ and\ \citenamefont
  {et~al.}(2017)}]{Zhang:2017gau}%
  \BibitemOpen
  \bibfield  {author} {\bibinfo {author} {\bibfnamefont {J.-H.}\ \bibnamefont
  {Zhang}}\ and\ \bibinfo {author} {\bibnamefont {et~al.}},\ }\href@noop {} {\
  (\bibinfo {year} {2017})},\ \Eprint {http://arxiv.org/abs/in preparation} {in
  preparation} \BibitemShut {NoStop}%
\bibitem [{\citenamefont {Ligeti}\ \emph {et~al.}(2008)\citenamefont {Ligeti},
  \citenamefont {Stewart},\ and\ \citenamefont {Tackmann}}]{Ligeti:2008ac}%
  \BibitemOpen
  \bibfield  {author} {\bibinfo {author} {\bibfnamefont {Z.}~\bibnamefont
  {Ligeti}}, \bibinfo {author} {\bibfnamefont {I.~W.}\ \bibnamefont {Stewart}},
  \ and\ \bibinfo {author} {\bibfnamefont {F.~J.}\ \bibnamefont {Tackmann}},\
  }\href@noop {} {\bibfield  {journal} {\bibinfo  {journal} {Phys. Rev.}\
  }\textbf {\bibinfo {volume} {D78}},\ \bibinfo {pages} {114014} (\bibinfo
  {year} {2008})},\ \Eprint {http://arxiv.org/abs/0807.1926} {arXiv:0807.1926
  [hep-ph]} \BibitemShut {NoStop}%
%%CITATION = ARXIV:0807.1926;%%
\end{thebibliography}%

\makeatother
\endgroup
